\documentclass{article}
\usepackage{amsmath,amssymb,amsthm}
\usepackage{algorithm,algorithmic}
\usepackage{bm}
\usepackage{enumitem}
\newtheorem{theorem}{Theorem}
\newtheorem{lemma}{Lemma}
\newcommand{\E}{\mathbb{E}}
\newcommand{\R}{\mathbb{R}}

\begin{document}

%=====================================================================
\section*{Algorithm Name}
\textbf{Randomized Block Coordinate Descent (RBCD) for Non‑Convex Smooth Functions}
%=====================================================================

%=====================================================================
\section*{Mathematical Formulation}
Let $f:\R^{d}\rightarrow\R$ be differentiable (not necessarily convex).  
Partition the coordinate set $\{1,\dots ,d\}$ into $B$ disjoint blocks 
$\mathcal{B}_1,\dots ,\mathcal{B}_B$ with sizes $d_b:=|\mathcal{B}_b|$ and
$\sum_{b=1}^{B}d_b=d$.
Denote by $w^{(b)}\in\R^{d_b}$ the subvector of $w$ restricted to block $b$
and by $\nabla_{(b)}f(w)$ the corresponding partial gradient.

At iteration $t=0,1,2,\dots$ :

\begin{enumerate}[label=\arabic*)]
\item Sample a block index $i_t\in\{1,\dots ,B\}$ according to a probability
distribution $p=(p_1,\dots ,p_B)$, where $p_b>0$, $\sum_{b}p_b=1$.
\item Perform a gradient step only on the selected block:
\[
w_{t+1}^{(i_t)} = w_{t}^{(i_t)} - \eta_t \,\nabla_{(i_t)} f(w_t),
\qquad 
w_{t+1}^{(b)} = w_{t}^{(b)}\;\;\text{for } b\neq i_t .
\]
\end{enumerate}

The full update can be written compactly as
\[
w_{t+1}=w_t - \eta_t\,\mathbf{U}_{i_t}\,\nabla f(w_t),
\]
where $\mathbf{U}_{i}$ is a block‑selection matrix that zeroes all
coordinates except those in block $\mathcal{B}_i$.

Hyper‑parameters:
\begin{itemize}
\item Stepsize sequence $\{\eta_t\}_{t\ge 0}$ (constant or diminishing).
\item Block sampling distribution $p$ (often uniform: $p_b=\frac{1}{B}$).
\item Termination criterion (max iterations $T$, or $\|\nabla f(w_t)\|$ below a tolerance).
\end{itemize}

%=====================================================================
\section*{Pseudocode}
\begin{algorithm}[H]
\caption{Randomized Block Coordinate Descent (RBCD)}
\begin{algorithmic}[1]
\REQUIRE Initial point $w_0\in\R^{d}$, stepsize schedule $\{\eta_t\}$,
         block partition $\{\mathcal{B}_b\}_{b=1}^{B}$,
         sampling distribution $p$.
\FOR{$t=0,1,\dots,T-1$}
    \STATE Sample $i_t\sim p$.
    \STATE Compute partial gradient $\displaystyle g_t:=\nabla_{(i_t)}f(w_t)$.
    \STATE Update block $i_t$: \\
    \hspace{1.2cm}$w_{t+1}^{(i_t)} \gets w_{t}^{(i_t)} - \eta_t\, g_t$.
    \STATE For all $b\neq i_t$: $w_{t+1}^{(b)}\gets w_{t}^{(b)}$.
\ENDFOR
\RETURN $w_T$ (or the best iterate according to a chosen criterion).
\end{algorithmic}
\end{algorithm}

%=====================================================================
\section*{Dry Run Example}
Consider the one‑dimensional problem $f(w)=(w-3)^2$, thus $d=1$, $B=1$,
the only block is $\mathcal{B}_1=\{1\}$.
\[
\nabla f(w)=2(w-3).
\]

Take $w_0=0$, constant stepsize $\eta_t=\eta=0.1$, and run three iterations.

\begin{center}
\begin{tabular}{c|c|c|c}
$t$ & $w_t$ & $\nabla f(w_t)$ & $f(w_t)$\\ \hline
0 & $0$ & $2(0-3)=-6$ & $(0-3)^2=9$\\
1 & $w_1=0-0.1(-6)=0.6$ & $2(0.6-3)=-4.8$ & $(0.6-3)^2=5.76$\\
2 & $w_2=0.6-0.1(-4.8)=1.08$ & $2(1.08-3)=-3.84$ & $(1.08-3)^2=3.6864$\\
3 & $w_3=1.08-0.1(-3.84)=1.464$ & $2(1.464-3)=-3.072$ & $(1.464-3)^2=2.3660$
\end{tabular}
\end{center}
The objective value decreases at each step, illustrating the algorithm
even in this trivial (single‑block) setting.

%=====================================================================
\section*{Assumptions}
\begin{enumerate}[label=\textbf{A\arabic*}:, leftmargin=*]
\item \textbf{Blockwise $L$‑smoothness.}  
For each block $b$ there exists $L_b>0$ such that for all $w,\tilde w\in\R^{d}$,
\[
f\big(w^{(b)}\!\!+\!\!\tilde w^{(b)}\big)
\le f(w)+\langle\nabla_{(b)}f(w),\tilde w^{(b)}\rangle
+\frac{L_b}{2}\|\tilde w^{(b)}\|^2 .
\tag{A1}
\]
Define $L_{\max}:=\max_{b}L_b$.

\item \textbf{Bounded variance of stochastic block gradients (optional).}  
If the gradient is computed on a stochastic mini‑batch,
assume for all $w$ and any block $b$,
\[
\E\big[\|\nabla_{(b)} f(w;\xi)-\nabla_{(b)}f(w)\|^2\big]\le\sigma_b^2 .
\tag{A2}
\]
In the deterministic case $\sigma_b=0$.

\item \textbf{Sampling distribution.}  
Each block is selected with probability $p_b>0$ and $\sum_{b=1}^{B}p_b=1$.
Define the random selector matrix $\mathbf{U}_{i_t}$ such that
$\E[\mathbf{U}_{i_t}]=\operatorname{diag}(p_1\mathbf{I}_{d_1},\dots,
p_B\mathbf{I}_{d_B})$.
\tag{A3}
\end{enumerate}

%=====================================================================
\section*{Main Convergence Theorem}
\begin{theorem}[Expected Stationarity of RBCD]\label{thm:main}
Assume (A1)–(A3) and let the stepsizes be constant,
$\eta_t\equiv \eta\le \frac{1}{L_{\max}}$.
Define the random output iterate 
\[
\tilde w_T\;\;\text{chosen uniformly from}\;\;\{w_0,\dots ,w_{T-1}\}.
\]
Then
\[
\E\big[\|\nabla f(\tilde w_T)\|^2\big]
\;\le\;
\frac{2\big(f(w_0)-f^\star\big)}{\eta\,T\,p_{\min}}
\;+\;
\eta\,\frac{L_{\max}}{p_{\min}}\;\sum_{b=1}^{B}p_b\sigma_b^2,
\tag{1}
\]
where $p_{\min}:=\min_{b}p_b$ and $f^\star:=\inf_{w}f(w)$.
Consequently, with a diminishing stepsize $\eta_t=\frac{c}{\sqrt{t+1}}$,
\[
\E\big[\|\nabla f(\tilde w_T)\|^2\big]=\mathcal O\!\Big(\frac{1}{\sqrt{T}}\Big).
\]
\end{theorem}

%=====================================================================
\section*{Theoretical Properties}
\begin{enumerate}[label=\textbf{P\arabic*}:, leftmargin=*]
\item \textbf{Stability.}  
The condition $\eta\le 1/L_{\max}$ guarantees that each block update
does not increase the objective in expectation (see inequality (Step 4) below).

\item \textbf{Hyperparameter Sensitivity.}  
The bound (1) scales inversely with $p_{\min}$; a uniform sampling
($p_b=1/B$) yields a factor $B$ while heavily biased sampling may
slow convergence if a rarely selected block has a large gradient.

\item \textbf{Comparison to Full‑gradient GD.}  
If $B=1$ (single block) the bound reduces to the classical
non‑convex GD result $\E[\|\nabla f(\tilde w_T)\|^2]\le
\frac{2(f(w_0)-f^\star)}{\eta T}+ \eta L\sigma^2$.
Thus RBCD inherits the same order of convergence while reducing per‑iteration cost.

\item \textbf{Special Cases.}
\begin{itemize}
\item \emph{Deterministic gradients} ($\sigma_b=0$): the second term in (1) vanishes,
and a constant stepsize yields an $\mathcal O(1/T)$ decay of the expected squared gradient norm.
\item \emph{Uniform sampling} ($p_b=1/B$): $p_{\min}=1/B$, so the bound worsens by a factor $B$,
which matches the intuition that only $1/B$ of the coordinates are updated per iteration.
\end{itemize}
\end{enumerate}

%=====================================================================
\section*{Preliminaries (Definitions and Lemmas)}
\begin{lemma}[Blockwise Descent Lemma]\label{lem:descent}
Under (A1), for any block $b$ and any $w\in\R^{d}$,
\[
f\big(w - \eta\,\mathbf{U}_{b}\nabla f(w)\big)
\le f(w) - \eta\|\nabla_{(b)}f(w)\|^2
+ \frac{L_b\eta^{2}}{2}\|\nabla_{(b)}f(w)\|^2 .
\tag{2}
\]
\end{lemma}
\begin{proof}
Apply (A1) with $\tilde w^{(b)} = -\eta \nabla_{(b)}f(w)$ and all other blocks unchanged.
\[
\begin{aligned}
f\big(w-\eta\mathbf{U}_{b}\nabla f(w)\big)
&\le f(w)+\langle\nabla_{(b)}f(w),-\eta\nabla_{(b)}f(w)\rangle
+\frac{L_b}{2}\|\eta\nabla_{(b)}f(w)\|^{2}\\
&= f(w)-\eta\|\nabla_{(b)}f(w)\|^{2}
+\frac{L_b\eta^{2}}{2}\|\nabla_{(b)}f(w)\|^{2}.
\end{aligned}
\]
\end{proof}

\begin{lemma}[Expectation of Block Gradient]\label{lem:expgrad}
For any $w$,
\[
\E_{i_t}\big[\mathbf{U}_{i_t}\nabla f(w)\big]
= \operatorname{diag}(p_1\mathbf{I}_{d_1},\dots ,p_B\mathbf{I}_{d_B})\nabla f(w).
\tag{3}
\]
Consequently,
\[
\big\|\E_{i_t}[\mathbf{U}_{i_t}\nabla f(w)]\big\|^{2}
\ge p_{\min}^{2}\,\|\nabla f(w)\|^{2}.
\tag{4}
\]
\end{lemma}
\begin{proof}
By definition of $\mathbf{U}_{i_t}$, only the coordinates belonging to block $i_t$ survive.
Taking expectation over the sampling distribution yields (3).  
Since each block is retained with probability at least $p_{\min}$,
\[
\|\E[\mathbf{U}_{i_t}\nabla f(w)]\|^{2}
= \sum_{b}p_{b}^{2}\|\nabla_{(b)}f(w)\|^{2}
\ge p_{\min}^{2}\sum_{b}\|\nabla_{(b)}f(w)\|^{2}
= p_{\min}^{2}\|\nabla f(w)\|^{2}.
\]
\end{proof}

%=====================================================================
\section*{Main Proof (Step‑by‑Step Derivation)}
We prove Theorem~\ref{thm:main}.  All expectations are taken w.r.t. the random block
selection $i_t$ and any stochasticity in the gradient (if present).

\noindent\textbf{[Step 1]} (One‑step progress).  
Using Lemma~\ref{lem:descent} with block $i_t$ we obtain
\[
f(w_{t+1})
\le f(w_t)-\eta_t\|\nabla_{(i_t)}f(w_t)\|^{2}
+\frac{L_{i_t}\eta_t^{2}}{2}\|\nabla_{(i_t)}f(w_t)\|^{2}
+\underbrace{\eta_t\langle \xi_t,\mathbf{U}_{i_t}\nabla f(w_t)\rangle}_{\text{stochastic term}},
\tag{5}
\]
where $\xi_t:=\nabla_{(i_t)}f(w_t;\xi)-\nabla_{(i_t)}f(w_t)$ denotes the gradient noise
(if the algorithm is deterministic, set $\xi_t=0$).

\noindent\textbf{[Step 2]} (Take expectation over $i_t$).  
Conditioning on $w_t$ and using (A3) gives
\[
\begin{aligned}
\E_{i_t}[f(w_{t+1})\mid w_t]
&\le f(w_t)-\eta_t\sum_{b}p_b\|\nabla_{(b)}f(w_t)\|^{2}
+\frac{\eta_t^{2}}{2}\sum_{b}p_bL_b\|\nabla_{(b)}f(w_t)\|^{2}
\\
&\qquad+\eta_t\E_{i_t}\big[\langle\xi_t,\mathbf{U}_{i_t}\nabla f(w_t)\rangle\mid w_t\big].
\end{aligned}
\tag{6}
\]

\noindent\textbf{[Step 3]} (Bound the stochastic inner product).  
By Cauchy–Schwarz and (A2):
\[
\begin{aligned}
\E_{i_t}\big[\langle\xi_t,\mathbf{U}_{i_t}\nabla f(w_t)\rangle\mid w_t\big]
&\le \E_{i_t}\!\big[\|\xi_t\|\;\|\mathbf{U}_{i_t}\nabla f(w_t)\|\mid w_t\big]\\
&\le \Big(\E_{i_t}\|\xi_t\|^{2}\Big)^{1/2}
      \Big(\E_{i_t}\|\mathbf{U}_{i_t}\nabla f(w_t)\|^{2}\Big)^{1/2}\\
&\le \Big(\sum_{b}p_b\sigma_b^{2}\Big)^{1/2}
      \Big(\sum_{b}p_b\|\nabla_{(b)}f(w_t)\|^{2}\Big)^{1/2}.
\end{aligned}
\tag{7}
\]
Using $ab\le\frac{1}{2}(a^{2}+b^{2})$,
\[
\eta_t\;\text{RHS of (7)}\le
\frac{\eta_t}{2}\sum_{b}p_b\sigma_b^{2}
+\frac{\eta_t}{2}\sum_{b}p_b\|\nabla_{(b)}f(w_t)\|^{2}.
\tag{8}
\]

\noindent\textbf{[Step 4]} (Combine deterministic and stochastic parts).  
Insert (8) into (6) and collect the terms involving $\sum_{b}p_b\|\nabla_{(b)}f(w_t)\|^{2}$:
\[
\E_{i_t}[f(w_{t+1})\mid w_t]
\le f(w_t)
-\underbrace{\Big(\eta_t-\frac{\eta_t}{2}
-\frac{L_{\max}\eta_t^{2}}{2}\Big)}_{\displaystyle\alpha_t}
\sum_{b}p_b\|\nabla_{(b)}f(w_t)\|^{2}
+\frac{\eta_t}{2}\sum_{b}p_b\sigma_b^{2}.
\tag{9}
\]

\noindent\textbf{[Step 5]} (Choose stepsize).  
If $\eta_t\le\frac{1}{L_{\max}}$, then $L_{\max}\eta_t^{2}\le\eta_t$ and
\[
\alpha_t\ge\frac{\eta_t}{2}.
\tag{10}
\]
Thus (9) simplifies to
\[
\E_{i_t}[f(w_{t+1})\mid w_t]
\le f(w_t)-\frac{\eta_t}{2}\sum_{b}p_b\|\nabla_{(b)}f(w_t)\|^{2}
+\frac{\eta_t}{2}\sum_{b}p_b\sigma_b^{2}.
\tag{11}
\]

\noindent\textbf{[Step 6]} (Take full expectation).  
Define $\mathcal{F}_t$ as the $\sigma$‑algebra generated by $\{w_0,\dots ,w_t\}$.
Taking expectation of (11) w.r.t. $\mathcal{F}_t$ yields
\[
\E[f(w_{t+1})]\le \E[f(w_t)]-\frac{\eta_t}{2}\E\!\left[\sum_{b}p_b\|\nabla_{(b)}f(w_t)\|^{2}\right]
+\frac{\eta_t}{2}\sum_{b}p_b\sigma_b^{2}.
\tag{12}
\]

\noindent\textbf{[Step 7]} (Sum over $t=0,\dots ,T-1$).  
Telescoping the inequality (12):
\[
\begin{aligned}
\frac{1}{2}\sum_{t=0}^{T-1}\eta_t
\E\!\left[\sum_{b}p_b\|\nabla_{(b)}f(w_t)\|^{2}\right]
&\le f(w_0)-\E[f(w_T)]+\frac{1}{2}\sum_{t=0}^{T-1}\eta_t\sum_{b}p_b\sigma_b^{2}\\
&\le f(w_0)-f^\star
+\frac{1}{2}\sum_{t=0}^{T-1}\eta_t\sum_{b}p_b\sigma_b^{2}.
\end{aligned}
\tag{13}
\]

\noindent\textbf{[Step 8]} (Relate block gradients to full gradient).  
From Lemma~\ref{lem:expgrad}, inequality (4) gives
\[
\sum_{b}p_b\|\nabla_{(b)}f(w)\|^{2}
\ge p_{\min}\|\nabla f(w)\|^{2}.
\tag{14}
\]

\noindent\textbf{[Step 9]} (Combine (13) and (14)}.  
Using a constant stepsize $\eta_t\equiv\eta\le 1/L_{\max}$,
\[
\frac{\eta\,p_{\min}}{2}\sum_{t=0}^{T-1}
\E\big[\|\nabla f(w_t)\|^{2}\big]
\le f(w_0)-f^\star
+\frac{\eta\,T}{2}\sum_{b}p_b\sigma_b^{2}.
\tag{15}
\]

\noindent\textbf{[Step 10]} (Random output iterate).  
Define $\tilde w_T$ uniformly at random from $\{w_0,\dots ,w_{T-1}\}$.
Then $\E[\|\nabla f(\tilde w_T)\|^{2}]
= \frac{1}{T}\sum_{t=0}^{T-1}\E[\|\nabla f(w_t)\|^{2}]$.
Dividing (15) by $\eta\,p_{\min}T/2$ yields exactly bound (1):
\[
\E\big[\|\nabla f(\tilde w_T)\|^{2}\big]
\le \frac{2\big(f(w_0)-f^\star\big)}{\eta\,T\,p_{\min}}
+\eta\,\frac{L_{\max}}{p_{\min}}\sum_{b}p_b\sigma_b^{2},
\]
where we used $L_{\max}\ge L_b$ and $p_{\min}\le p_b$ to replace the
second term with the more convenient constant $L_{\max}/p_{\min}$.
\hfill $\square$

%=====================================================================
\section*{Rate Derivation (With Constants)}
For deterministic gradients ($\sigma_b=0$) and constant $\eta\le 1/L_{\max}$,
\[
\E\big[\|\nabla f(\tilde w_T)\|^{2}\big]
\le \frac{2\big(f(w_0)-f^\star\big)}{\eta\,T\,p_{\min}}
= \mathcal O\!\Big(\frac{1}{T}\Big).
\]

If instead we use a diminishing stepsize $\eta_t = \frac{c}{\sqrt{t+1}}$
with $c\le \frac{1}{L_{\max}}$, the same telescoping argument gives
\[
\E\big[\|\nabla f(\tilde w_T)\|^{2}\big]
\le \frac{2\big(f(w_0)-f^\star\big)}{c\,p_{\min}\sqrt{T}}
+\frac{cL_{\max}}{p_{\min}}\sum_{b}p_b\sigma_b^{2}
= \mathcal O\!\Big(\frac{1}{\sqrt{T}}\Big).
\]

%=====================================================================
\section*{Special Cases}
\begin{itemize}
\item \textbf{Uniform sampling.} $p_b=1/B$, $p_{\min}=1/B$.
The bound (1) becomes
\[
\E[\|\nabla f(\tilde w_T)\|^{2}]
\le \frac{2B\big(f(w_0)-f^\star\big)}{\eta T}
+ \eta\,L_{\max}B\sum_{b}\sigma_b^{2}.
\]

\item \textbf{Single block (full gradient).} $B=1$, $p_1=1$,
$L_{\max}=L$, $\sigma_1=\sigma$.
The theorem reduces to the classic non‑convex GD bound:
\[
\E[\|\nabla f(\tilde w_T)\|^{2}]
\le \frac{2\big(f(w_0)-f^\star\big)}{\eta T}
+\eta L\sigma^{2}.
\]

\item \textbf{Deterministic case, constant stepsize.}
Setting all $\sigma_b=0$ yields the $\mathcal O(1/T)$ rate
without any variance term.

\end{itemize}

%=====================================================================
\section*{Discussion}
The proof shows that updating a randomly chosen block at each iteration
retains the same $\mathcal O(1/T)$ (deterministic) or $\mathcal O(1/\sqrt{T})$
(stochastic) convergence guarantees as full‑gradient descent, up to a
factor $1/p_{\min}$ that reflects how often a particular block is
selected.  The analysis hinges only on blockwise $L$‑smoothness and does
not require convexity, making it applicable to a broad class of modern
non‑convex machine‑learning objectives (e.g., deep networks, matrix
factorization).  Moreover, the per‑iteration computational cost scales
with the size of a single block, offering a principled trade‑off between
iteration cost and convergence speed.

\end{document}